\uuid{noB2}
\titre{Convergence}
\auteur{bodin ; hanani ; mzari}
\organisation{exo7}
\licence{CC BY-NC-SA 4.0}
\datecreate{2020-01-20}
\niveau{L1}
\module{Analyse}
\chapitre{Suite}
\sousChapitre{Autre}
\difficulte{2}
\theme{Suites, convergence}
\exo7qcmid{lille-1.suites.020}

\contenu{
\begin{question}
\qclassification{121.00}
\qtags{motcle=convergence}
Soit $\displaystyle u_n=\sum _{k=1}^n\frac{1}{k(k+1)}$. Quelles sont les bonnes réponses ?
\begin{answers}
    \bad{La suite $(u_n)$ est croissante non majorée.}
    \bad{La suite $(u_n)$ est divergente.}
    \good{Pour tout $n\geq 1$, $\displaystyle u_n=1-\frac{1}{n+1}$.}
    \good{$(u_n)$ est convergente et $\displaystyle \lim _{n\to +\infty}u_n=1$.}
\end{answers}
\begin{explanations}
Le terme $u_n$ est une somme télescopique. En effet, on vérifie que, pour tout $k\geq 1$,
$$\frac{1}{k(k+1)}=\frac{1}{k}-\frac{1}{k+1}\Rightarrow u_n=1-\frac{1}{n+1}.$$
Donc $(u_n)$ est convergente et sa même limite est $1$.
\end{explanations}
\end{question}
}
